% ============================================================
%  Kapitel 1: Allgemein -- TilVista und MainWindow
% ============================================================
\chapter{Allgemein: TilVista und das MainWindow}

\section{Projektueberblick}

TilVista ist eine Qt6-Desktop-Anwendung zur Verwaltung und
Anzeige von Mediendateien -- Bildern und Videos.
Der Name ist ein Kunstwort aus \emph{tilvelde}
(norw. zufaellig) und \emph{vista} (span./ital. Aussicht).
Version \version{0.5.42} ist die aktuelle Entwicklungsversion.

Das Programm gliedert sich in drei Haupt-Werkzeuge,
die als Tabs im Hauptfenster angeordnet sind:

\begin{itemize}[noitemsep]
  \item \textbf{AleaVue} -- Zufaellige Vollbild-Diashow aus
    einem Verzeichnis (Bilder, S-Taste zum Bookmarken)
  \item \textbf{SattumaPic} -- Zufaelliger Dateiaufruf mit
    Vorschau (beliebige Dateitypen)
  \item \textbf{MadoLudus} -- Konfigurations-Tab, oeffnet ein
    eigenstaendiges \klasse{MadoludusWindow} fuer Bilder und
    Videos mit Flip-Book-Modus (\textbf{neu ab v0.5.42})
\end{itemize}

Zwei persistente Datenbanken erganzen die Werkzeuge:
\begin{itemize}[noitemsep]
  \item \textbf{kaivo (DB1)} -- Verzeichnis-Datenbank mit
    gecachten Dateilisten (\datei{.dshow} / \datei{.catalogue})
  \item \textbf{shujuko (DB2)} -- Lesezeichen fuer einzelne
    Dateien, je kaivo-Eintrag eine eigene JSON-Datei
\end{itemize}

\section{Architekturueberblick}

Das Programm folgt einer flachen, signalgesteuerten Hierarchie.
\klasse{MainWindow} besitzt alle Tabs und gemeinsam genutzten
Komponenten. Tabs kommunizieren \emph{ausschliesslich} ueber
Signale und Slots mit \klasse{MainWindow} -- niemals direkt
untereinander.

\begin{center}
\ttfamily\small
\begin{tabular}{|c|}
\hline
\textbf{MainWindow (QMainWindow)} \\
\hline
DirBar -- globale Verzeichnisauswahl (oben) \\
\hline
QTabWidget \\
\quad AleaVueTab \textbar{} SattumaPicTab \textbar{}
MadoludusTab \textbar{} ShortcutsTab \\
\hline
StatusBar + SecretIndicator (unten) \\
\hline
\end{tabular}
\end{center}

\begin{neuin}
\textbf{v0.5.42-Refactor Madoludus:}
Der fruehre \klasse{MadolodosTab} war ein
Monolith (Konfiguration + Wiedergabe in einer Klasse).
Jetzt trennt \klasse{MadoludusTab} die Konfiguration vom
Wiedergabe-Fenster (\klasse{MadoludusWindow}), das beim
Start-Klick als eigenstaendiges \klasse{QMainWindow}
geoeffnet wird. Das Fenster ist rahmenlos, immer im
Vordergrund und per Maus verschiebbar.

Neu ebenfalls: \verb|closeEvent| in \klasse{MainWindow}
ruft \methode{flushPendingWrites} auf beiden Datenbanken auf,
damit laufende Hintergrund-Saves beim Beenden nicht
abgebrochen werden.
\end{neuin}

\klasse{ShujukoPanel} ist eine Besonderheit:
Es wird in \klasse{MainWindow} erstellt, aber als geteilter
Rohzeiger (shared, not owned) sowohl an
\klasse{AleaVueTab} als auch an \klasse{SattumaPicTab}
weitergereicht. Der Eigentuemer ist \klasse{MainWindow}.

\section{Dateien: mainwindow.h}

\begin{lstlisting}[style=header, caption={mainwindow.h -- vollstaendige Deklaration}]
#pragma once
#include <QMainWindow>
#include <QStringList>

// Forward-Deklarationen: kein #include noetig,
// solange nur Zeiger deklariert werden.
class DirBar;
class AleaVueTab;
class DirDatabasePanel;
class MadoludusTab;     // v0.5.42: war MadolodosTab
class ShortcutsTab;
class ShujukoPanel;
class SattumaPicTab;
class QCloseEvent;      // v0.5.42: neu (fuer closeEvent)
class QLabel;
class QShortcut;
class QTabWidget;

class MainWindow : public QMainWindow
{
    Q_OBJECT
public:
    explicit MainWindow(QWidget* parent = nullptr);

protected:
    // v0.5.42: wartet auf laufende Hintergrund-Saves
    void closeEvent(QCloseEvent* event) override;

private slots:
    void onDirChanged(const QString& path);
    void onDirFromDb(const QString& path);
    void onDb1FilesLoaded(const QString& path,
                          const QStringList& imageFiles,
                          const QStringList& allFiles);
    void onActiveEntryChanged(const QString& entryName,
                               const QString& sourceDir);
    void toggleSecretMode();

private:
    void updateSecretIndicator();

    DirBar*        m_dirBar        = nullptr;
    QTabWidget*    m_tabs          = nullptr;
    ShujukoPanel*  m_shujuko       = nullptr;
    AleaVueTab*    m_aleaVueTab    = nullptr;
    SattumaPicTab* m_sattumaPicTab = nullptr;
    MadoludusTab*  m_madoludusTab  = nullptr;
    ShortcutsTab*  m_shortcutsTab  = nullptr;

    QLabel*    m_secretIndicator = nullptr;
    QShortcut* m_secretShortcut  = nullptr;
    bool       m_secretMode      = false;
};
\end{lstlisting}

\subsubsection{C++-Analyse}

\paragraph{\#pragma once}
Ein nicht-standardisierter, aber universell unterstuetzter
Include-Guard. Statt
\verb|#ifndef MAINWINDOW_H / #define / #endif|
sorgt \verb|#pragma once| dafuer, dass die Datei pro
Kompiliereinheit nur einmal eingelesen wird -- kuerzer und
fehlerunanfaelliger.

\paragraph{Forward-Deklarationen}
Statt \verb|#include "aleavuetab.h"| reicht im Header
\verb|class AleaVueTab;| -- solange nur Zeiger auf den Typ
deklariert werden, muss der Compiler die Klasse noch nicht
vollstaendig kennen. Weniger Abhaengigkeiten, schnellere
Kompilierung (Pimpl-Idiom light).

\paragraph{Mitglieder-Initialisierung mit = nullptr}
In C++11 und spaeter kann man Pointer-Member direkt im
Klassenkoerper initialisieren. Das verhindert
uninitialisierte Zeiger (undefined behavior) und erspart
explizite Konstruktorinitialisierer.

\subsubsection{Qt-Analyse}

\paragraph{QMainWindow}
Die Basisklasse liefert Menueleiste, Toolbar-Bereich,
CentralWidget und Statusleiste -- vorgefertigt, kein
manueller Aufbau noetig. \klasse{MainWindow} erbt alles
durch \verb|: public QMainWindow|.

\paragraph{Q\_OBJECT}
Pflicht-Makro fuer jede Klasse mit Signalen oder Slots.
Der Meta Object Compiler (\texttt{moc}) generiert daraus
den C++-Code, der das Signal-Slot-System implementiert.
\cmake{AUTOMOC ON} in \datei{SetupQt.cmake} erledigt den
\texttt{moc}-Aufruf automatisch beim Build.

\paragraph{private slots}
Slots koennen jede Sichtbarkeit haben. \verb|private slots|
bedeutet: Die Slots koennen nur intern oder ueber
\verb|connect()| genutzt werden, nicht von aussen als
normale Methode aufgerufen.

\section{Dateien: mainwindow.cpp}

\begin{lstlisting}[caption={Konstruktor: Fenster aufbauen}]
MainWindow::MainWindow(QWidget* parent)
    : QMainWindow(parent)
{
    setWindowTitle(
        QString("TilVista  \u2013  %1")
        .arg(TV_APPVERSION_DISPLAY));
    setMinimumSize(900, 640);

    // Icon aus Qt-Ressourcen (resources.qrc)
    const QIcon appIcon(":/icons/tilvista.ico");
    if (!appIcon.isNull()) setWindowIcon(appIcon);

    // Zentrales Widget + vertikales Layout
    auto* central = new QWidget;
    setCentralWidget(central);
    auto* mv = new QVBoxLayout(central);
    mv->setContentsMargins(8,8,8,8);
    mv->setSpacing(4);

    m_dirBar = new DirBar;
    connect(m_dirBar, &DirBar::directoryChanged,
            this, &MainWindow::onDirChanged);
    mv->addWidget(m_dirBar);

    // ShujukoPanel: Eigentuemer MainWindow, geteilt mit Tabs
    m_shujuko = new ShujukoPanel(this);

    m_tabs = new QTabWidget;

    // Lambda gibt den aktuellen Verzeichnis-Zustand weiter
    m_aleaVueTab = new AleaVueTab(
        [this]{ return m_dirBar->directory(); }, m_shujuko);
    m_sattumaPicTab = new SattumaPicTab(
        [this]{ return m_dirBar->directory(); }, m_shujuko);
    m_madoludusTab  = new MadoludusTab(
        [this]{ return m_dirBar->directory(); });
    m_shortcutsTab  = new ShortcutsTab;

    // Signalkette: DB1 -> MainWindow -> alle betroffenen Tabs
    DirDatabasePanel* db1 = m_aleaVueTab->dbPanel();
    connect(m_aleaVueTab, &AleaVueTab::requestDirChange,
            this, &MainWindow::onDirFromDb);
    connect(db1, &DirDatabasePanel::dirLoaded,
            this, &MainWindow::onDb1FilesLoaded);
    connect(db1, &DirDatabasePanel::activeEntryChanged,
            this, &MainWindow::onActiveEntryChanged);
    connect(db1, &DirDatabasePanel::requestShujukoValidation,
            m_shujuko, &ShujukoPanel::validateFiles);

    m_tabs->addTab(m_aleaVueTab,    "\U0001f5bc  AleaVue");
    m_tabs->addTab(m_sattumaPicTab, "\U0001f3b2  SattumaPic");
    m_tabs->addTab(m_madoludusTab,  "\U0001f39e  MadoLudus");
    m_tabs->addTab(m_shortcutsTab,  "About");
    mv->addWidget(m_tabs);

    // Schloss-Symbol in der Statusleiste
    m_secretIndicator = new QLabel("\U0001f512");
    statusBar()->addPermanentWidget(m_secretIndicator);

    m_secretShortcut = new QShortcut(
        QKeySequence("Ctrl+Alt+F8"), this);
    m_secretShortcut->setContext(Qt::ApplicationShortcut);
    connect(m_secretShortcut, &QShortcut::activated,
            this, &MainWindow::toggleSecretMode);
}
\end{lstlisting}

\subsubsection{C++-Analyse: Lambda als Callback}

\begin{lstlisting}
m_aleaVueTab = new AleaVueTab(
    [this]{ return m_dirBar->directory(); }, m_shujuko);
\end{lstlisting}

Ein \textbf{Lambda-Ausdruck} wird als
\verb|std::function<QString()>| uebergeben:
\begin{itemize}[noitemsep]
  \item \verb|[this]| -- der Lambda-Closure haelt eine
    Kopie des \texttt{this}-Zeigers von \klasse{MainWindow}
  \item \verb|{ return m_dirBar->directory(); }| -- liefert
    beim Aufruf das aktuelle Verzeichnis aus der DirBar
\end{itemize}
Damit muss \klasse{AleaVueTab} kein \klasse{MainWindow}
kennen -- \emph{Dependency Inversion Principle}.

\subsubsection{v0.5.42 neu: closeEvent}

\begin{lstlisting}[caption={closeEvent: warten auf laufende Saves}]
void MainWindow::closeEvent(QCloseEvent* event) {
    // Blockiert kurz (Event-Loop pumpen), bis in-flight
    // Saves auf Disk abgeschlossen sind.
    // Details: DirDatabasePanel::flushPendingWrites().
    m_aleaVueTab->dbPanel()->flushPendingWrites();
    m_shujuko->flushPendingWrites();
    QMainWindow::closeEvent(event);  // Basisklasse!
}
\end{lstlisting}

Ohne diese Methode koennte ein "Save to DB"-Thread,
der genau beim Beenden des Programms lief, mitten im
Schreiben abgebrochen werden -- das JSON waere dann
korrupt.

\subsubsection{Qt-Analyse: Signalkette}

Wenn der Nutzer einen DB-Eintrag doppelklickt:

\begin{enumerate}[noitemsep]
  \item \klasse{DirDatabasePanel} emittiert
    \signal{dirLoaded(path, imageFiles, allFiles)}
  \item \slot{MainWindow::onDb1FilesLoaded} empfaengt
  \item Leitet \verb|path + allFiles| an
    \klasse{SattumaPicTab} weiter
  \item Leitet \verb|path + allFiles| an
    \klasse{MadoludusTab} weiter
\end{enumerate}

Alle Tabs bekommen die Daten -- ohne sich gegenseitig
zu kennen.

\section{Gemeinsam genutzte Hilfsdateien}

\subsection{pathutils.h / pathutils.cpp}

Alle anwendungsweiten Pfade, Dateinamen-Hilfsfunktionen
und plattformspezifische Calls sind im Namespace \ns{TV}
gebuendelt.

\begin{lstlisting}[style=header, caption={pathutils.h}]
namespace TV {
    QString scriptDir();    // Verzeichnis der .exe
    QString kaivoDir();     // .../optegnelser/kaivo/
    QString kaivoDbPath();  // .../kaivo/tilvista_dirs.json
    QString shujukoPath(const QString& kaivoEntryName);
    QString logDir();       // .../optegnelser/
    QString safeFilename(const QString& name);
    QString autoEntryName(const QString& dirPath);

    const QStringList& imageSuffixes(); // .jpg, .png, ...
    const QStringList& videoSuffixes(); // .mp4, .mkv, ...
    QStringList allMediaSuffixes();

    void openPath(const QString& path);
    void selectInExplorer(const QString& path);
    void preventSleep();  // Win32: SetThreadExecutionState
    void restoreSleep();
}
\end{lstlisting}

\paragraph{Statische lokale Variable -- Singleton-Initialisierung}
\begin{lstlisting}
const QStringList& imageSuffixes() {
    static const QStringList s = {
        ".jpg", ".jpeg", ".png", ".webp",
        ".heic", ".heif", ".avif", ".tif", ...
    };
    return s;
}
\end{lstlisting}
\verb|static| bei einer lokalen Variable: Initialisierung
genau einmal beim ersten Aufruf, danach im Speicher bis
Programmende. Rueckgabe als \verb|const QStringList&|
(Referenz) -- keine Kopie bei jedem Aufruf.

\paragraph{Raw-String-Literal R"(...)" und safeFilename}
\begin{lstlisting}
QString safeFilename(const QString& name) {
    static const QString bad = R"(\/:*?"<>|)";
    QString out;
    for (QChar c : name)
        out += bad.contains(c) ? QChar('_') : c;
    while (out.startsWith('.')) out.remove(0,1);
    return out.isEmpty()
        ? QStringLiteral("entry") : out;
}
\end{lstlisting}
\verb|R"(...)"| ist ein \textbf{Raw-String-Literal} (C++11):
Backslashes und Anfuehrungszeichen brauchen kein Escape.
\verb|QStringLiteral| erzeugt den String zur Compile-Zeit.

\paragraph{Plattformspezifischer Code mit Q\_OS\_WIN}
\begin{lstlisting}
void preventSleep() {
#ifdef Q_OS_WIN
    SetThreadExecutionState(
        ES_CONTINUOUS | ES_SYSTEM_REQUIRED);
#endif
}
\end{lstlisting}
Qt-eigene Praeprozessormakros (\verb|Q_OS_WIN|,
\verb|Q_OS_MAC|, \verb|Q_OS_LINUX|) machen den Code
portable: Auf Nicht-Windows-Systemen kompiliert die
Funktion als leere Huelle.

\subsection{config.h / config.cpp}

\klasse{Config} liest beim Programmstart einmalig
\datei{tilvista\_config.json} und stellt Einstellungen
als statische Klassenmitglieder bereit.

\begin{lstlisting}[style=header, caption={config.h}]
class Config {
public:
    static void     load();
    static QVariant get(const QString& key);
    static bool     getBool(const QString& key);
private:
    static QJsonObject s_data;
    static QJsonObject s_defaults;
};
\end{lstlisting}

\paragraph{Statische Membervariablen -- Klassenweites Singleton}
\verb|static QJsonObject s_data| gehoert der Klasse,
nicht einer Instanz. Es gibt genau ein \verb|s_data|
im gesamten Programm -- klassisches Singleton-Muster
ohne Overhead.

\begin{qtinfo}
\textbf{QJsonObject als Konfigurationsspeicher:}
Qt stellt \klasse{QJsonObject}, \klasse{QJsonArray} und
\klasse{QJsonDocument} fuer vollstaendiges JSON-Handling
bereit. Parsing:
\begin{lstlisting}
QJsonParseError err;
auto doc = QJsonDocument::fromJson(f.readAll(), &err);
if (err.error == QJsonParseError::NoError
    && doc.isObject())
    s_data = doc.object();
\end{lstlisting}
Dokumentation: \url{https://doc.qt.io/qt-6/qjsondocument.html}
\end{qtinfo}
